% !TeX program = xelatex
% 中文起步稿：Overleaf 设置 XeLaTeX，Main document 选择 starter.tex。
\documentclass[10pt,chinese]{njulink}

\njulinksetup{
  type={研究报告},
  id={NJU-LINK / RESEARCH NOTES},
  short-title={NJU-LINK 研究报告},
  subtitle={用清晰的结构连接问题、方法与证据},
  keywords={研究问题；方法设计；实验评估}
}
\title{NJU-LINK 研究报告}
\author[1]{作者姓名}
\author[1]{合作者姓名}
\affiliation[1]{NJU-LINK}
\metadata[日期]{\today}
% \contribution[*]{共同第一作者}
% \correspondence{\email{name@example.org}}

\abstract{在此概述研究问题、核心方法、主要证据与适用范围。
摘要应当能够独立阅读，并与正文中的实际结果保持一致。
这是一份可直接替换内容的起步模板；请在完成研究后填写具体结论。}

\begin{document}
\maketitle

\section{研究背景}
说明研究场景、已有方法的局限，以及本文需要回答的具体问题。
模板支持中文、English、数学符号 $\mathcal{D}$ 与公式的混合排版。

\begin{njubox}[title={研究要点}]
用一至两句话写清楚研究问题，或总结一条已由实验支持的发现。
\end{njubox}

\section{方法}
先定义输入、输出和假设，再介绍方法的主要步骤。
例如，对 $N$ 个样本的得分 $s_i$，平均分定义为
\begin{equation}
  \overline{s}=\frac{1}{N}\sum_{i=1}^{N}s_i.
\end{equation}

\section{实验与讨论}
描述数据、基线、评价指标和不确定性，并说明结论的适用范围。
图、表、引用、附录与参考文献的完整示例见 \texttt{main.tex}。

% 需要参考文献时，取消下列注释并在正文中添加 \citep{...} 或 \citet{...}。
% \bibliographystyle{assets/plainnat}
% \bibliography{references}
\end{document}
